\section{Proof Sketch}

This section contains proof sketches of Theorem \ref{lemma:approx}, \ref{thm:generalization} and Lemma \ref{lemma:gaussian}. 

\subsection{Proof Sketch of Theorem \ref{lemma:approx}}\label{sec:proof-sketch-approx}

We provide a proof sketch of Theorem \ref{lemma:approx} in this part and defer technical details of the proof to Appendix \ref{appendix:approx}.
The ReLU neural network in Theorem \ref{lemma:approx} is constructed in the following 5 steps:
\begin{enumerate}%[itemsep=1 pt,topsep=1 pt,parsep=1 pt]
    \item Choose region $S \subset \cD_{\rm data}$ only on which we use ReLU neural networks to approximate $f^*$.
    \item Construct a covering of $S$ with hypercubes and then divide these hypercubes into several groups, so that neural networks constructed with respect to each group have nonoverlapping supports.
    \item Implement ReLU neural networks to assign given input and estimated function value to corresponding hypercube.
    \item Approximate $f^*$ by a Taylor polynomial and then implement a ReLU neural network to approximate Taylor polynomial on each hypercube.
    \item Sum up all the sub-neural-networks and take maximum to approximate $f^*$.
\end{enumerate}

\textbf{Step 1. Space separation.} Firstly, we divide $\cD_{\rm data}$ into some region $S \subset \cD_{\rm data}$ with high probability measure and $S^c = \cD_{\rm data} $ with large volume. By Assumption \ref{assumption:pdata_cover}, for any sufficiently small $r,\tau>0$ and some constant $c_0>1$, there exists $S \subset \cD_{\rm data}$ such that $N_r(S) \leq c_0 N(r;\tau) \leq c_0 r^{-p}$ for some positive constant $p=p(r,\tau)$ and $P_{\rm data}(S^c) \leq \tau$. Intuitively, we only need to approximate $f^*$ on $S$ while $S^c$ is negligible due to its small probability measure.  Therefore, in the following steps, we only design a covering for $S$ and approximate $f^*$ in each hypercube of the covering. 

\textbf{Step 2. Grouping hypercubes.} Let $\cC$ be a minimum set of hypercubes with side length $r$ covering $S$. Then we partition $\cC$ into $\cC_1,\ldots,\cC_J$ such that each subset $\cC_j$ is composed of hypercubes separated by $r$ from each other. Lemma \ref{lem:covering_division} shows that the number of $\cC_j$'s is at most a constant dependent of $d$.

As a consequence, we group hypercubes into several subsets of $\cC$ so that  constructed neural networks with respect to each hypercube in $\cC_j$ have nonoverlapping support. 


\textbf{Step 3. Hypercube Determination.} 
This step is to assign the given input $x$ and estimated function value $y$ to the hypercube where they belong. To do so, we design a neural network to approximate function $(x,y) \mapsto y\ind_I(x)$ where $I \in \cC$ is some hypercube. To make functions positive, we firstly consider approximating $f_0 = f^* + 2$. Notice that $f_0 \in \cH(\beta,\cD_{\rm data}, 3)$ and $1 \leq f_0(x) \leq 3 $ for any $x \in \cD_{\rm data}$. 

For any fixed $I \in \cC$, we define the center of $I$ as $(\iota_1, \dots, \iota_d)$.
Then we construct a neural network $g_I^{\text{ind},r}: \cD_{\rm data} \times \RR_\geq \to \RR_\geq$ with the form:
\begin{align}
    g_I^{\text{ind},r} (x,y) =  4 \relu \bigg(\sum_{i=1}^d \hat{\ind}_{I, i}^r(x_i) + \frac{y}4 - d \bigg),\label{eq:identical}
\end{align}
where $\hat{\ind}^r_{I, i}: \RR \to [0, 1]$ is the approximated indicator function given by
\begin{equation*}
  \hat{\ind}^r_{I, i}(z) = \begin{cases}
  %0 & \text{if } z \leq \iota_i - r,\\
  \frac{z - (\iota_i - r)}{r/2} & \text{if } \iota_i - r < z \leq \iota_i - \frac{r}{2},\\
  1 & \text{if } \iota_i - \frac{r}{2} < z \leq \iota_i + \frac{r}{2},\\
  \frac{(\iota_i + r) - z}{r/2} & \text{if } \iota_i + \frac{r}{2} < z \leq \iota_i + r,\\
  0 & \text{otherwise. } %\iota_i+r < z.
\end{cases}
\end{equation*}
We claim that neural network $g_I^{\text{ind},r}$ approximates function $(x,y) \mapsto y\ind_I(x)$. Moreover, Appendix \ref{sec:realize-ind} provides the explicit realization of $g_I^{\text{ind},r}$ by selecting specific weight matrices and intercepts.

 
\textbf{Step 4. Taylor Approximation.} 
In each cube $I \in \cC$, we locally approximate $f^*$ by a Taylor polynomial of degree $\lfloor\beta\rfloor$ and then we define a neural network to approximate this Taylor polynomial. Firstly, we cite the following lemma to evaluate the difference between any $\beta$-H\"{o}lder function and its Taylor polynomial:
\begin{lemma}[Lemma A.8 in \citet{PETERSEN2018296}]\label{lemma: taylor_approx}
  Fix any $f \in \cH(\beta, \cD_{\rm data})$ with $\norm{f}_{\cH(\beta, \cD_{\rm data})} \leq 1$ and $\bar{x} \in S$.
  Let $\bar{f}(x)$ be the Taylor polynomial of degree $\lfloor\beta\rfloor$ of $f$ around $\bar{x}$, namely,
  \begin{equation*}
    \bar{f}(x) = \sum_{ |\alpha| \leq \lfloor\beta\rfloor} \frac{\partial^\alpha f(\bar{x})}{\alpha!} (x - \bar{x})^\alpha.
  \end{equation*}
  Then, $|f(x) - \bar{f}(x)| \leq d^\beta \norm{x - \bar{x}}^\beta$ holds for any $x \in \cD_{\rm data} $.%, where $\underline{C}$ is a constant depending on $\beta, D$.
\end{lemma}

Next, we design an $m$-dimensional multiple output neural network $g^{\poly}_\epsilon=(g^{\poly}_{\epsilon,1},\dots,g^{\poly}_{\epsilon,m})$ to estimate multiple Taylor polynomials in each output. The existence of such neural network is ensured in the following lemma, which is a straightforward extension of Lemma 18 in \citet{nakada2020adaptive}.

\begin{lemma}[Taylor approximation on S]\label{lemma:poly}
  Fix any $m \in \NN_+$. Let $\{c_{k, \alpha}\} \subset [-1, 1]$ for $1 \leq k \leq m$. Let $\{x_k\}_{k=1}^m \subset S$.
  Then there exist $c_1^{\poly} = c_1^{\poly}(\beta, d, p)$, $c_2^{\poly} = c_2^{\poly}(\beta, d, p)$ and $s^{\poly}_1 = s^{\poly}_1(\beta, d, p)$  such that for any sufficiently small $\epsilon>0$, there is a neural network $g^{\poly}_\epsilon$ which satisfies the followings:
  \begin{enumerate}%[itemsep=1.3pt,topsep=1.3pt,parsep=1.2pt]
    \item $\sup_{x \in S} \bigl| g^{\poly}_{\epsilon,k}(x) - \sum_{|\alpha| < \beta} c_{k, \alpha} (x - x_k)^\alpha \bigr| \leq \epsilon$ for any $k=1,\ldots,m$,
    \item $L(g^{\poly}_\epsilon) \leq 1 + (2 + {\log_2\beta})(11 + (1+\beta)/p)$,
    \item $B(g^{\poly}_\epsilon) \leq c_1^{\poly} R_S^{\beta s_1^{\poly}}\epsilon^{-s_1^{\poly}}$,
    \item $K(g^{\poly}_\epsilon) \leq c_2^{\poly} (R_S^p \epsilon^{-p/\beta} + m)$.
  \end{enumerate}
\end{lemma}

For any cube $I_k \in \cC$, we take $f_{I_k}(x)$ as a Taylor polynomial function as with setting $\bar{x} \leftarrow x_{I_k}$ and $f \leftarrow f_0$ in Lemma \ref{lemma: taylor_approx}. Then we define a neural network to approximate $f_{I_k}$, which is an $\epsilon/2$-accuracy Taylor polynomial of $f_0$. Let $g^{\poly}_{\epsilon/2}$ be a neural network constructed in Lemma \ref{lemma:poly} with $\epsilon \leftarrow \epsilon/2$, $m \leftarrow N_r(S)$, $(x_k)_{k=1}^{m} \leftarrow (x_{I_k})_{k=1}^{N_r(S)}$, and
$(c_{k, \alpha})_{k=1}^{m} \leftarrow (\partial^\alpha f(x_{I_k})/\alpha!)_{k=1}^{N_r(S)}$ appearing in Lemma \ref{lemma:poly}.
Then, we obtain
\begin{align}
  \sup_{k=1,\dots,N_r(S)} \sup_{x \in S} \bigl| f_{I_k}(x) - g^{\poly}_{\epsilon/2,k}(x) \bigr| \leq \frac{\epsilon}{2}. \label{eq:poly}
\end{align}

In addition, we construct a neural network to aggregate the outputs of $g^{\poly}_{\epsilon/2}$. Define a neural network $g^{\text{filter}}_k : \RR^{d + N_r(S)} \to \RR^{d+1}$ which picks up the first $d$ inputs and $(d + k)$-th input as
\begin{equation*}
    g^{\text{filter}}_k (z)= \begin{pmatrix}\II_d & e_k^\top \\ \bf{0}_d & e_k^\top \end{pmatrix}z,~\text{ for } k=1,\dots,N_r(S).
\end{equation*}
Then we design a neural network $g^{\text{simul}}_{\epsilon/2}: \RR^d \to \RR^{N_r(S)}$ that simultaneously estimates Taylor polynomial at each cube. Specifically, $g^{\text{simul}}_{\epsilon/2}$ is formulated as below  
\begin{equation}\label{eq:simul}
\begin{split}
    g^{\text{simul}}_{\epsilon/2} = \bigl(g_{I_1}^{\text{ind},r}\circ g^{\text{filter}}_1 ,\ldots,g_{I_{ N_r(S)}}^{\text{ind},r}
    \circ g^{\text{filter}}_{ N_r(S)} \bigr) \circ (g^{Id}_{d,L},  g^{\poly}_{\epsilon/2}),
\end{split}
\end{equation}
where $g^{Id}_{d,L}: \RR^d \to \RR^d$ is the neural network version of the identity function whose number of layers is equal to $L(g^{\poly}_{\epsilon/2})$.


\textbf{Step 5. Construction of Neural Networks.} 
In this step, we construct a neural network $g^{f_0}_\epsilon$ to approximate $f_0 = f^* + 2 $. Let $g^{\max,5^d}$ be the neural network version of the maximize function over $5^d$ numbers. Besides, define
\begin{equation*}
    g^{\text{sum}}(z_1, \dots, z_{N_r(S)}) = \biggl( \sum_{I_k \in \cC_1} z_k, \dots,  \sum_{I_k \in \cC_{N_r(S)}} z_k \biggr),
\end{equation*}
which aims to sum up the output of $g^{\text{simul}}_{\epsilon/2}$ in each subset of covering $\cC_j$. 

Now we are ready to define $g^{f_0}_\epsilon$. Let $g^{f_0}_\epsilon:=g^{\max,5^d} \circ g^{\text{sum}} \circ g^{\text{simul}}_{\epsilon/2}$. Equivalently,  $g^{f_0}_\epsilon$ can be written as  $g^{f_0}_\epsilon=\max_{j \in [5^d]} \sum_{I_k \in \cC_j} g^{\text{simul}}_{\epsilon/2,k}$. Then we come to bound the approximation error of $g^{f_0}_\epsilon$. When $x \in S$, there exists some $I \in \cC$ such that $x \in I$.  Based on the method to construct neural networks, we have
\begin{align*}
    g^{f_0}_\epsilon(x) = \max_{I_k \in \neig(I)}  g^{\text{simul}}_{\epsilon/2,k}(x)  \leq\max_{I_k \in \neig(I)}  g^{\poly}_{\epsilon/2,k}(x),
\end{align*}
where $\neig(I)=\{ I' \in \cC | (I \oplus 3r/2)\cap I' \neq \emptyset \}$ denotes the $3r/2$-neighborhood of hypercube $I$.
In other words, when computing $g^{f_0}_\epsilon(x)$, we only need to take maximum over the estimated function value within hypercubes near $x$.
%For a further parameter tuning, we set $\gamma = D^{-1}\qty(3 M)^{-1/\beta} \varepsilon^{1/\beta}$.

Given sufficiently small $\epsilon > 0$, 
the error is bounded as
\begin{align*}
    | g^{f_0}_\epsilon (x) - f_0(x) | 
    \leq& \max_{I_k \in \neig(I)}  \bigl| g^{\poly}_{\epsilon/2,k}(x)-f_0(x) \bigr| \\
    \leq& \max_{I_k \in \neig(I)}  \bigl| g^{\poly}_{\epsilon/2,k}(x)-f_{I_k}(x) \bigr|+ \max_{I_k \in \neig(I)}  \bigl| f_{I_k}(x) -f_0(x) \bigr| \\
    \leq& \frac{\varepsilon}{2} + d^\beta \biggl(\frac{3r}{2} \biggr)^\beta \leq \epsilon,
\end{align*}
where the last inequality follows from \eqref{eq:poly} and Lemma \ref{lemma: taylor_approx}. Detailed derivation of approximation error is deferred to Appendix \ref{sec:approx-error-bound}. In terms of parameter tuning, we choose $r = d^{-1} \epsilon^{1/\beta}/2$. 

To extend results of $f_0$ to $f^*$, we implement a neural network $g^{\text{mod}}(z)=(-z+1)\circ\relu(-z+2)\circ \relu(z-1)$ and consider $g^{f^*}_\epsilon=g^{\text{mod}}\circ g^{f_0}_\epsilon$ to obtain the desired approximation error $\epsilon$ for any $x\in S$. Then we evaluate the approximation error with respect to $L^2$-norm:
\begin{align*}
    \Bigl\| g^{f^*}_\epsilon-f^* \Bigr\|_{L^2(P_{\rm data})} 
    =\biggl(\int_{S} +\int_{S^c} \biggr) \Bigl(g^{f^*}_\epsilon(x)-f^*(x) \Bigr)^2 \ud P_{\rm data}(x) \leq \epsilon^2 + 4\tau. %\norm{f^*}_\infty^2 .
\end{align*}
This follows from the aforementioned approximation error within $S$, boundedness of $f^*$ and neural networks, as well as the property that out-of-$S$ probability is upper bounded, i.e. $P_{\rm data}(S^c) \leq \tau$.

Finally, we sum up sizes of all the sub-neural-networks and thus obtain the network size of $g^{f^*}_\epsilon$. See Appendix \ref{sec:parameter} for detailed calculation.%Details of the proof for Theorem \ref{lemma:approx} are deferred to Appendix \ref{appendix:approx}.




\subsection{Proof Sketch of Theorem \ref{thm:generalization}}\label{proof:sketch-generalization}

Proof of Theorem \ref{thm:generalization} follows a standard statistical decomposition, i.e. decomposing the mean squared error of estimator $\hat{f}$ into a squared bias term and a variance term. We bound the bias and variance separately, where the bias is tackled by the approximation results given in Theorem \ref{lemma:approx} and the variance is bounded using the metric entropy arguments. 
%\citep{van1996weak, gyorfi2006distribution}. 
Details of the proof for Theorem\ref{thm:generalization} are provided in Appendix \ref{sec:proof-generalization}.
At first, we decompose the $L_2$ risk as follows:
\begin{align*}
\EE \norm{\hat{f}-f^*}^2_{L^2(P_{\rm data})} = \underbrace{2 \EE \biggl [\frac{1}{n}\sum_{i=1}^n (\hat{f}(x_i) - f^*(x_i))^2 \biggr]}_{T_1}  +\underbrace{\EE\norm{\hat{f}-f^*}^2_{L^2(P_{\rm data})} - 2 \EE\biggl [\frac{1}{n}\sum_{i=1}^n (\hat{f}(x_i) - f^*(x_i))^2 \biggr]}_{T_2},
\end{align*}
where $T_1$ reflects the squared bias of using neural networks to estimate $f^*$ and $T_2$ is the variance term.
% The proof of Theorem \ref{thm:stat} contains two parts for the estimates of $T_1$ and $T_2$, respectively. 

\textbf{Step 1. Bounding bias term $\bf{T_1}$.}
Since $T_1$ is the empirical $L_2$ risk of $\hat{f}$ evaluated on the samples $\{x_i\}_{i=1}^n$, we relate $T_1$ to the empirical risk by rewriting $f^*(x_i) = y_i - \xi_i$,  so that we can apply the approximation error to bound the minimal empirical risk achieved by $\hat{f}$. After some basic calculation, we have
\begin{align*}
T_1 %=& 2 \EE \left[\frac{1}{n} \sum_{i=1}^n (\hat{f}(x_i) - y_i + \xi_i)^2 \right] \nonumber \\
%\overset{(i)}{=}& 2 \EE \left[\frac{1}{n} \sum_{i=1}^n \left[(\hat{f}(x_i) - y_i)^2 + 2 \xi_i \hat{f}(x_i) - \xi_i^2 \right]\right] \nonumber \\
 %=& 2 \EE \biggl[\inf_{f \in \cF(R, L, B, K)} \frac{1}{n}\sum_{i=1}^n \Bigl[(f(x_i) - y_i)^2 \\
 %&\quad - \xi_i^2 + 2 \xi_i \hat{f}_n(\xb_i) \Bigr] \biggr] \nonumber \\
 %\overset{(ii)}
 \leq  2 \inf_{f \in \cF( L, B, K)} \norm{ f(x) - f^*(x)}^2_{L^2(P_{\rm data})}  + 4 \EE \left[\frac{1}{n} \sum_{i=1}^n \xi_i \hat{f}(x_i) \right].
\end{align*}
%Equality $(i)$ is obtained by expanding the square, where the cross term $\EE[\xi_i y_i] = \EE[\xi_i (f_0(\xb_i) + \xi_i)] = \EE[\xi_i^2]$ due to the independence between $\xb_i$ and $\xi_i$. Inequality $(ii)$ invokes the Jensen's inequalty to pass the expectation. To obtain term $(A)$, we expand $(f(\xb_i) - y_i)^2 = (f(\xb_i) - f_0(\xb_i) - \xi_i)^2$, and observe the cancellation of $-\xi_i^2$. 
Note that the first term is the squared approximation error of neural networks, which can be controlled by Theorem \ref{lemma:approx}. We bound the second term by quantifying the complexity of the network class $\cF( L, B, K)$. A precise upper bound of $T_1$ is given in the following lemma.
%whose proof follows a similar argument in \citet[Lemma 4]{schmidt2017nonparametric}.
\begin{lemma}\label{lemma:t1}
Fix the neural network class $\cF( L, B, K)$. For any $\delta \in (0, 1)$, there exists some constant $c>0$, such that
\begin{align*}
T_1 \leq &~ c \inf_{f \in \cF( L, B, K)} \norm{ f(x) - f^*(x)}^2_{L^2(P_{\rm data})} + c \sigma^2 \frac{\log \cN_2(\delta, \cF( L, B, K)) + 2}{n} \\
&\quad+  c\bigg(   \sqrt{\frac{\log \cN_2(\delta, \cF(L, B, K)) + 2}{n}} + 1 \bigg)\sigma \delta,
\end{align*}
where $\cN_2(\delta, \cF(L, B, K))$ denotes the $\delta$-covering number of $\cF(L, B, K)$ with respect to the $L^2$ norm, i.e., there exists a discretization of $\cF(L, B, K)$ into $\cN_2(\delta, \cF(L, B, K))$ distinct elements, such that for any $f \in \cF$, there is $\bar{f}$ in the discretization satisfying $\norm{\bar{f} - f}_2 \leq \epsilon$.
\end{lemma}

\textbf{Step 2. Bounding variance term $\bf{T_2}$.}
We observe that $T_2$ is the difference between the population risk of $\hat{f}$ and its empirical risk. However, bounding this difference is distinct from traditional concentration results because of the scaling factor $2$ before the empirical risk. To do this, we divide the empirical risk into two parts and use a higher-order moment (fourth moment) to bound one part. Using a Bernstein-type inequality, we are able to show that $T_2$ converges at a rate of $1/n$, and the upper bound for this is shown in the following lemma.
\begin{lemma}\label{lemma:t2}
For any $\delta \in (0, 1)$, there exists some constant $c'>0$, such that
\begin{align*}
T_2 \leq& \frac{c'}{3n} \log \cN_2(\delta/4H, \cF(L, B, K))+ c'\delta.
\end{align*}
\end{lemma}



\textbf{Step 3. Covering number of neural networks.}
The upper bounds of $T_1$ and $T_2$ in Lemma \ref{lemma:t1} and \ref{lemma:t2} both rely on the covering number of the network class $\cF(R, \kappa, L, p, K)$. In this step, we present an upper bound for the covering number $\cN_2(\delta, \cF(L, B, K) )$ for a given a resolution $\delta > 0$. %Since each weight parameter in the network is bounded by a constant $\kappa$, we construct a covering by partitioning the range of each weight parameter into a uniform grid. By choosing a proper grid size, we show the following lemma.
\begin{lemma}[Covering number bound for $\cF$ (Lemma 21 in \citet{nakada2020adaptive})]\label{lemma:covering-number}
  Given $\delta>0$, the $\delta$-covering number of the neural network class $\cF(L,B,K)$ satisfies
  \begin{align*}
    \log \cN_2(\delta, \cF( L, B,K)) \leq K \log\biggl(\! \frac{2^L \sqrt{dL} K^{L/2} B^L R_S }{\sqrt{\delta^2-4\tau}}\! \biggr).
  \end{align*}
\end{lemma}

\textbf{Step 4. Bias-Variance Trade-off.}
Now we are ready to finish the proof of Theorem \ref{thm:generalization}. Combining the upper bounds of $T_1$ in Lemma \ref{lemma:t1} and $T_2$ in Lemma \ref{lemma:t2} together and substituting the covering number in  Lemma \ref{lemma:covering-number}, we obtain
\begin{align*}
\EE \norm{ \hat{f} - f^*}^2_{L^2(P_{\rm data})} 
        =  O \!\biggl(\! \tau + d^{2\beta} r^{2\beta} \!+\sigma\delta+  \frac{\sigma^2 K}{n} \log\biggl( \frac{ \sqrt{dL} K^{L/2} B^L R_S }{\sqrt{\delta^2-4\tau}}  \biggr)\!\biggr),
\end{align*}
where we set approximation error to be $d^{2\beta} r^{2\beta}$. 
Plug in our choice of $(L,B,K)$, and choose $\delta = r^{2\beta}$. Then we can conclude
\begin{align*}
    \EE \norm{ \hat{f} - f^*}^2_{L^2(P_{\rm data})} =  O\biggl(   \tau +  \sigma r^{2\beta} 
    +\frac{\sigma^2}{n} \biggl(\frac{R_S}{r} \biggr)^p \log\biggl( \frac{(R_S/r)^p}{ r^{4\beta}-4\tau}\biggr) \biggr).
\end{align*}




\subsection{Proof Sketch of Lemma \ref{lemma:gaussian}}\label{sec:gaussian}

In this section, we present our basic idea to construct $S(R,r;p)$
and the proof sketch of Lemma \ref{lemma:gaussian}. For simplicity of proof, we assume $Q=I$ so that $\Sigma=\Lambda=\diag(\lambda_1^2,\ldots,\lambda_d^2)$. The detailed proof is given in Appendix \ref{proof:gaussian}, which can be easily extended to the case when $Q$ is not an identity matrix. The proof of Theorem \ref{thm:generalization-exp} is given in Appendix \ref{sec:dim-gaussian}.

Given the Gaussian sample distribution, we hope to choose some region in $S \subset \RR^d$ with high probability measure and effective Minkowski dimension $p<d$. Then we can only apply neural networks to approximate $f^*$ within each cube of the small covering of $S$ and thereby significantly reduce the network size.
In literature, it is common to truncate the ambient space within a hyper-ellipsoid for Gaussian distribution \citep{ellis2007multivariate,pakman2014exact}. Similarly, we consider the 'thick' low-dimensional hyper-ellipsoid $S(R,r;p)$ defined in \eqref{eq:set}. 
%\begin{align*}
%     S(R,r;p) := \bigg\{z \in \RR^d \bigg|& \sum_{i=1}^p \frac{z_i^2}{\lambda_i^2} \leq R^2, \notag \\
%    &|z_j| \leq \frac{r}{2} \text{ for } j=p+1,\ldots,d  \bigg\}. \label{eq:set2}
%\end{align*}
%where $R,r>0$ and $p\in \NN_+$ are independent parameters. 
Then we construct a minimal cover of $S(R,r;p)$ as a union of nonoverlapping hypercubes with side length $r$, which is equal to the thickness of $S(R,r;p)$. In particular, this cover contains multiple layers of hypercubes to cover the first $p$ dimensions of $S(R,r;p)$ while only needs one layer for the rest dimensions. Intuitively, we only learn about hypercubes that cover the first $p$ dimensions without paying extra effort to study dimensions with low probability density. 


A natural question arising from the construction of $S(R,r;p)$ is how to select a proper dimension $p$. To address this problem, we first notice that each side length of the $p$-dimension hyper-ellipsoid is supposed to be greater than the side length of hypercubes $r$, i.e. $2\lambda_i R \geq r$ for $i=1,\ldots,p$, so that we would not waste hypercubes to cover dimensions with too small side length. For simplicity of calculation, we choose $p>0$ that satisfies $\lambda_p^{-1} = 2R/r$ for any given $R,r>0$.  

Now we come to prove Lemma \ref{lemma:gaussian}. Firstly, we compute the probability outside $S(R,r;p)$. By union bound, this probability can be upper bounded by two parts, the probability out of hyper-ellipsoid for the first $p$ dimensions and the probability out of hypercube with side length $r$ for the rest $d-p$ dimensions. The first part is equal to the tail bound of $p$-dimensional standard Gaussian by the construction of $S(R,r;p)$. The second part can be solved similarly by linearly transforming each dimension to be standard Gaussian.

Then we calculate the covering number of $S(R,r;p)$. Notice that the first $p$ dimensions of $S(R,r;p)$ is contained in a $p$-dimensional hyper-rectangle with side length $2\lambda_i R$ for $i=1,\ldots,p$, while only one hypercube is required to cover the $j$-th dimension for $j=p+1,\ldots,d$. Therefore, the $r$-covering number can be upper bounded by $\prod_{i=1}^p (2\lambda_i R)/r^p$.
